% ============================================================
%  5. Discussion & Reflection
% ============================================================
\section{Discussion \& Reflection}
\label{sec:discussion}

\subsection{Design Principles Revisited}

回顾从V5到V9CA\_AB的迭代，有三条经验最值得保留：

\begin{enumerate}[leftmargin=*, itemsep=4pt]
    \item \textbf{受控增量修正比强调制更稳定}：在本任务中，$h' = h + \text{offset}$比FiLM式乘性调制更可靠；训练策略上，pairwise $\to$ BCE $\to$ DirectionalReg的渐进课程也优于一步到位。
    \item \textbf{训练信号比模型容量更关键}：MLP和GRU的弱表现，以及T3\_NoS1/T2\_NoCur的明显退化，都说明高噪声金融数据中不能只依赖更复杂的模型。
    \item \textbf{逐项验证是必要工程纪律}：V11一次堆叠三项未经验证的改动后明显退化，说明本类项目中每个改动都需要独立验证。
\end{enumerate}

\subsection{What We Would Do Differently}

在项目时间约束下，仍有几个方向未能充分探索：

\begin{itemize}[leftmargin=*, itemsep=3pt]
    \item \textbf{Gate机制需要系统性验证}：E1\_NoEnc说明encoder整体重要，但目前还无法将收益退化干净归因到gate、peer dynamics或market context中的某一项。后续应按行业/市值分组统计gate值，并做更细粒度消融。
    \item \textbf{执行层仍需建模}：第~\ref{sec:trading}节的同花顺模拟交易显示，即使离线排序信号较强，真实模拟收益仍可能被盘中执行时点、成交价格、先卖后买的资金释放约束、仓位集中和人工择时显著削弱。尤其是本文离线协议采用每日top-20调仓，本身会带来较高换手；如果没有明确的换手惩罚、成交价假设和自动执行规则，交易层误差会直接侵蚀信号收益。后续应将predictor、scorer/selector和execution拆成独立模块，并固定调仓时点、仓位上限、换手约束和自动化下单协议。
    \item \textbf{新闻文本数据未能利用}：本次作业提供了2019年以来的每日新闻快讯数据。考虑到文本多模态处理的技术复杂度和项目时间约束，我们选择了聚焦于结构化量价特征——但新闻中的政策信号、行业事件等信息显然是peer dynamics和市场环境的重要补充。
\end{itemize}

\subsection{Conclusion}

本文提出了CLIME，一种面向A股日频top-K选股的Transformer排序模型。实验结果显示，两阶段训练、方向感知损失和受控市场信息注入共同提升了holdout收益和稳定性；内部分析也表明模型学到的并非纯噪声，而是以横截面排名、基本面和同业上下文为主的紧凑排序信号。

本工作仍有明确边界。第一，超参数搜索不充分，学习率、层数、head数等配置尚未系统扫描。第二，特征工程规模有限，67维特征虽覆盖主要因子类别，但尚未做大规模特征筛选。第三，CLIME仍是预测模型，不是完整交易系统；交易成本、流动性、涨跌停和执行时点都会削弱可实现收益。第四，同花顺模拟交易最终收益为$-19.40\%$，这暴露了离线alpha模型到真实交易之间的执行层缺口：日频排序信号若缺少固定调仓协议、自动化执行和仓位约束，可能被盘中价格路径和人工干预显著稀释。

更坦诚地说，CLIME当前解决的是一个相对“干净”的问题：在每天收盘后的信息集合下，预测下一交易日哪些股票更可能获得相对更好的收益。这个设定适合课程任务中的横向模型比较，也适合检验模型是否学到了可排序的alpha信号；但它并不直接解决真实操作中最麻烦的部分。实际交易中，首先要卖出旧仓才能买入新仓，卖出时点和买入时点并不天然同步；同时，A股日内K线波动明显，早盘上涨、午后回落或高开低走都会使“预测对了方向”和“账户赚到钱”之间出现差距。再加上top-20日调仓本身带来较高换手，短短约十个交易日的模拟比赛窗口又容易受到阶段性大盘环境影响，因此一次短期模拟交易的收益路径可能显著偏离100个交易日holdout上的平均信号表现。

这并不是说模拟交易结果可以被忽略。相反，它指出了本文最现实的不足：我们证明了模型在受控协议下能给出较强的次日排序信号，但还没有把这个信号完整转化为一个低换手、可成交、可复现的交易系统。若后续继续推进，最重要的改进不是再把predictor堆得更大，而是在模型分数之后加入换手惩罚、持仓惯性、日内成交规则、T+1约束、交易成本估计和风险预算。只有这些模块被形式化之后，离线top-20收益才更可能稳定转化为真实模拟收益。

因此，CLIME在本课程项目设定下是一个表现较强的次日排序模型，但若要走向实用量化系统，还需要在特征筛选、超参数搜索、组合构造和交易执行层上继续工程化。更准确地说，本文完成的是alpha predictor和信号解释，而不是完整自动交易系统；模拟交易中的失败案例恰好说明了这种模块边界的重要性。
